\chapter{Lp Spaces}
\label{ch:iii06}

$L^p$ spaces turn integral size into geometry. After identifying functions equal almost everywhere, the $L^p$ quantities become genuine norms for $p\ge1$, and the resulting spaces are complete.

\section*{Learning goals}
After this chapter, the reader should be able to:
\begin{itemize}
\item compute \(L^p\) norms and determine membership in \(L^p\) for explicit singular or decaying functions;
\item distinguish equality almost everywhere from equality as pointwise functions in \(L^p\) spaces;
\item prove completeness of \(L^p\) in the chapter's setting using Cauchy-sequence control;
\item compare \(L^p\) spaces on finite- and infinite-measure domains with explicit inclusions or counterexamples;
\item use truncation or simple-function approximation to construct dense subclasses;
\item analyze limiting behavior in \(L^p\) without confusing norm convergence with pointwise convergence;
\end{itemize}
\section*{Conceptual roadmap}
\[
\boxed{\text{structure}\;\longrightarrow\;\text{estimate}\;\longrightarrow\;\text{limit or representation}\;\longrightarrow\;\text{application}.}
\]

\section{Almost-everywhere classes}
Two measurable functions equal almost everywhere represent the same $L^p$ element.

\section{$L^p$ norms}
For $1\le p<\infty$, use $\|f\|_p=(\int|f|^p)^{1/p}$; for $p=\infty$, use essential supremum.


% BEGIN VOL03-EXPANSION III06-example-01
\begin{example}[Critical exponent for a singularity]\label{ex:iii06-expand-01}
On \((0,1)\), let \(f(x)=x^{-\alpha}\). Then
\[
\|f\|_p^p=\int_0^1x^{-\alpha p}dx,
\]
which is finite exactly when \(\alpha p<1\). The local singularity therefore
sets a sharp membership threshold across the \(L^p\) scale.
\end{example}
% END VOL03-EXPANSION III06-example-01
\section{Essential supremum}
The $L^\infty$ norm ignores exceptional values on null sets and records the least almost-everywhere uniform bound.

\section{Vector-space structure}
Hölder and Minkowski control sums and scalar multiples, making $L^p$ a normed vector space.

\section{Completeness}
Every $L^p$ Cauchy sequence converges in $L^p$ for $1\le p\le\infty$.

\section{Finite-measure inclusions}
On finite-measure spaces, $L^q\subset L^p$ when $q>p$, with an explicit measure factor.


% BEGIN VOL03-EXPANSION III06-example-02
\begin{example}[Finite measure gives Lq inside Lp]\label{ex:iii06-expand-02}
If \(\mu(X)<\infty\) and \(q>p\), Holder applied to \(|f|^p\cdot1\) gives
\[
\|f\|_p\le \mu(X)^{1/p-1/q}\|f\|_q.
\]
On a probability space the measure factor is one. The finite-measure
hypothesis is essential; this inclusion can fail on infinite-measure spaces.
\end{example}
% END VOL03-EXPANSION III06-example-02
\section{Convergence in measure}
For finite $p$, $L^p$ convergence implies convergence in measure by Chebyshev's inequality.

\section{Approximation}
Truncation, finite-measure cutoffs, and simple functions provide dense approximations in $L^p$ for $p<\infty$.


% BEGIN VOL03-EXPANSION III06-example-03
\begin{example}[Truncation approximates in Lp]\label{ex:iii06-expand-03}
For \(f\in L^p\), \(p<\infty\), let \(f_n=f1_{\{|f|\le n\}}\). Then
\[
|f-f_n|^p=|f|^p1_{\{|f|>n\}}\to0
\]
pointwise and is dominated by \(|f|^p\in L^1\). DCT yields
\(\|f-f_n\|_p\to0\). Thus unbounded \(L^p\) functions can first be reduced to
bounded approximants.
\end{example}
% END VOL03-EXPANSION III06-example-03
\section{Core structural results}

\begin{theorem}[Norm-zero criterion]\label{thm:iii06-01}
For $1\le p\le\infty$, $\|f\|_p=0$ iff $f=0$ almost everywhere.
\begin{proof}
For finite $p$, $\int|f|^p=0$ forces $|f|^p=0$ a.e.; for $p=\infty$, essential supremum zero means $|f|=0$ a.e.
\end{proof}
\end{theorem}

\begin{theorem}[Finite-measure inclusion]\label{thm:iii06-02}
If $\mu(X)<\infty$ and $1\le p<q\le\infty$, then $\|f\|_p\le\mu(X)^{1/p-1/q}\|f\|_q$.
\begin{proof}
For finite $q$, apply Hölder to $|f|^p\cdot1$ with exponent $q/p$ and its conjugate; the $q=\infty$ case follows directly.
\end{proof}
\end{theorem}

\begin{theorem}[Completeness of $L^p$]\label{thm:iii06-03}
For $1\le p\le\infty$, $L^p$ is complete.
\begin{proof}
Extract a subsequence with summable norm increments, use the resulting pointwise absolutely summable majorant to construct an a.e. limit, prove subsequence norm convergence, and recover the full Cauchy sequence.
\end{proof}
\end{theorem}

\section{Worked examples}

\begin{example}[Almost-everywhere classes diagnostic]\label{ex:iii06-worked-01}
Give a rigorous worked analysis of Almost-everywhere classes in the setting of this chapter. State the decisive definition or hypothesis and derive the central conclusion. Two measurable functions equal almost everywhere represent the same $L^p$ element. The decisive point is to use the chapter definition at the correct countable, norm, transform, or distributional level rather than rely on pointwise intuition alone.
\end{example}

\begin{example}[$L^p$ norms diagnostic]\label{ex:iii06-worked-02}
Give a rigorous worked analysis of $L^p$ norms in the setting of this chapter. State the decisive definition or hypothesis and derive the central conclusion. For $1\le p<\infty$, use $\|f\|_p=(\int|f|^p)^{1/p}$; for $p=\infty$, use essential supremum. The decisive point is to use the chapter definition at the correct countable, norm, transform, or distributional level rather than rely on pointwise intuition alone.
\end{example}

\begin{example}[Essential supremum diagnostic]\label{ex:iii06-worked-03}
Give a rigorous worked analysis of Essential supremum in the setting of this chapter. State the decisive definition or hypothesis and derive the central conclusion. The $L^\infty$ norm ignores exceptional values on null sets and records the least almost-everywhere uniform bound. The decisive point is to use the chapter definition at the correct countable, norm, transform, or distributional level rather than rely on pointwise intuition alone.
\end{example}

\begin{example}[Vector-space structure diagnostic]\label{ex:iii06-worked-04}
Give a rigorous worked analysis of Vector-space structure in the setting of this chapter. State the decisive definition or hypothesis and derive the central conclusion. Hölder and Minkowski control sums and scalar multiples, making $L^p$ a normed vector space. The decisive point is to use the chapter definition at the correct countable, norm, transform, or distributional level rather than rely on pointwise intuition alone.
\end{example}

\section{Solved dossiers}
Each dossier is a canonical solved problem. The provenance ledger records which retained corpus rules guided its topic and which dossiers were added freshly to complete the chapter.

\begin{problem}[Almost-everywhere classes diagnostic]\label{prob:iii06-dossier-01}
Give a rigorous worked analysis of Almost-everywhere classes in the setting of this chapter. State the decisive definition or hypothesis and derive the central conclusion.
\end{problem}
\begin{solution}
Two measurable functions equal almost everywhere represent the same $L^p$ element. The decisive point is to use the chapter definition at the correct countable, norm, transform, or distributional level rather than rely on pointwise intuition alone.
\end{solution}

\begin{problem}[$L^p$ norms diagnostic]\label{prob:iii06-dossier-02}
Give a rigorous worked analysis of $L^p$ norms in the setting of this chapter. State the decisive definition or hypothesis and derive the central conclusion.
\end{problem}
\begin{solution}
For $1\le p<\infty$, use $\|f\|_p=(\int|f|^p)^{1/p}$; for $p=\infty$, use essential supremum. The decisive point is to use the chapter definition at the correct countable, norm, transform, or distributional level rather than rely on pointwise intuition alone.
\end{solution}

\begin{problem}[Essential supremum diagnostic]\label{prob:iii06-dossier-03}
Give a rigorous worked analysis of Essential supremum in the setting of this chapter. State the decisive definition or hypothesis and derive the central conclusion.
\end{problem}
\begin{solution}
The $L^\infty$ norm ignores exceptional values on null sets and records the least almost-everywhere uniform bound. The decisive point is to use the chapter definition at the correct countable, norm, transform, or distributional level rather than rely on pointwise intuition alone.
\end{solution}

\begin{problem}[Vector-space structure diagnostic]\label{prob:iii06-dossier-04}
Give a rigorous worked analysis of Vector-space structure in the setting of this chapter. State the decisive definition or hypothesis and derive the central conclusion.
\end{problem}
\begin{solution}
Hölder and Minkowski control sums and scalar multiples, making $L^p$ a normed vector space. The decisive point is to use the chapter definition at the correct countable, norm, transform, or distributional level rather than rely on pointwise intuition alone.
\end{solution}

\begin{problem}[Completeness diagnostic]\label{prob:iii06-dossier-05}
Give a rigorous worked analysis of Completeness in the setting of this chapter. State the decisive definition or hypothesis and derive the central conclusion.
\end{problem}
\begin{solution}
Every $L^p$ Cauchy sequence converges in $L^p$ for $1\le p\le\infty$. The decisive point is to use the chapter definition at the correct countable, norm, transform, or distributional level rather than rely on pointwise intuition alone.
\end{solution}

\begin{problem}[Finite-measure inclusions diagnostic]\label{prob:iii06-dossier-06}
Give a rigorous worked analysis of Finite-measure inclusions in the setting of this chapter. State the decisive definition or hypothesis and derive the central conclusion.
\end{problem}
\begin{solution}
On finite-measure spaces, $L^q\subset L^p$ when $q>p$, with an explicit measure factor. The decisive point is to use the chapter definition at the correct countable, norm, transform, or distributional level rather than rely on pointwise intuition alone.
\end{solution}

\begin{problem}[Convergence in measure diagnostic]\label{prob:iii06-dossier-07}
Give a rigorous worked analysis of Convergence in measure in the setting of this chapter. State the decisive definition or hypothesis and derive the central conclusion.
\end{problem}
\begin{solution}
For finite $p$, $L^p$ convergence implies convergence in measure by Chebyshev's inequality. The decisive point is to use the chapter definition at the correct countable, norm, transform, or distributional level rather than rely on pointwise intuition alone.
\end{solution}

\begin{problem}[Approximation diagnostic]\label{prob:iii06-dossier-08}
Give a rigorous worked analysis of Approximation in the setting of this chapter. State the decisive definition or hypothesis and derive the central conclusion.
\end{problem}
\begin{solution}
Truncation, finite-measure cutoffs, and simple functions provide dense approximations in $L^p$ for $p<\infty$. The decisive point is to use the chapter definition at the correct countable, norm, transform, or distributional level rather than rely on pointwise intuition alone.
\end{solution}

\begin{problem}[Norm-zero criterion proof dossier]\label{prob:iii06-dossier-09}
Prove the following structural statement and identify the step where its main hypothesis is used: For $1\le p\le\infty$, $\|f\|_p=0$ iff $f=0$ almost everywhere.
\end{problem}
\begin{solution}
For finite $p$, $\int|f|^p=0$ forces $|f|^p=0$ a.e.; for $p=\infty$, essential supremum zero means $|f|=0$ a.e. This proof also explains why the stated hypothesis is structural rather than cosmetic.
\end{solution}

\begin{problem}[Finite-measure inclusion proof dossier]\label{prob:iii06-dossier-10}
Prove the following structural statement and identify the step where its main hypothesis is used: If $\mu(X)<\infty$ and $1\le p<q\le\infty$, then $\|f\|_p\le\mu(X)^{1/p-1/q}\|f\|_q$.
\end{problem}
\begin{solution}
For finite $q$, apply Hölder to $|f|^p\cdot1$ with exponent $q/p$ and its conjugate; the $q=\infty$ case follows directly. This proof also explains why the stated hypothesis is structural rather than cosmetic.
\end{solution}

\begin{problem}[Completeness of $L^p$ proof dossier]\label{prob:iii06-dossier-11}
Prove the following structural statement and identify the step where its main hypothesis is used: For $1\le p\le\infty$, $L^p$ is complete.
\end{problem}
\begin{solution}
Extract a subsequence with summable norm increments, use the resulting pointwise absolutely summable majorant to construct an a.e. limit, prove subsequence norm convergence, and recover the full Cauchy sequence. This proof also explains why the stated hypothesis is structural rather than cosmetic.
\end{solution}

\begin{problem}[Chapter synthesis]\label{prob:iii06-dossier-12}
Connect the main definitions and structural theorems of this chapter into one proof strategy. Explain what object is controlled, what estimate or closure principle is used, and what conclusion becomes available.
\end{problem}
\begin{solution}
$L^p$ spaces turn integral size into geometry. After identifying functions equal almost everywhere, the $L^p$ quantities become genuine norms for $p\ge1$, and the resulting spaces are complete. The recurring strategy is to encode the object in the chapter's natural measurable, normed, Fourier, or distributional structure, establish the controlling estimate, and then pass to the desired limit or representation.
\end{solution}

\section{Exercises with complete solutions}
The shorter exercise layer remains separate from the solved dossiers.

\begin{exercise}\label{exr:iii06-01}
State and apply the central principle behind Almost-everywhere classes. Give one immediate consequence in the setting of this chapter.
\end{exercise}
\begin{hint}
Compute \(\int |f|^p\) and examine the singularity and tail separately; both regions must be integrable.
\end{hint}
\begin{solution}
Two measurable functions equal almost everywhere represent the same $L^p$ element. As an immediate consequence, the same principle can be inserted into later convergence, approximation, transform, or weak-form arguments whenever its hypotheses are verified.
\end{solution}

\begin{exercise}\label{exr:iii06-02}
State and apply the central principle behind $L^p$ norms. Give one immediate consequence in the setting of this chapter.
\end{exercise}
\begin{hint}
Remember that \(L^p\) elements are equivalence classes: changing a function on a null set does not change its class.
\end{hint}
\begin{solution}
For $1\le p<\infty$, use $\|f\|_p=(\int|f|^p)^{1/p}$; for $p=\infty$, use essential supremum. As an immediate consequence, the same principle can be inserted into later convergence, approximation, transform, or weak-form arguments whenever its hypotheses are verified.
\end{solution}

\begin{exercise}\label{exr:iii06-03}
State and apply the central principle behind Essential supremum. Give one immediate consequence in the setting of this chapter.
\end{exercise}
\begin{hint}
For a Cauchy sequence in \(L^p\), extract a rapidly Cauchy subsequence and use summable norm differences to build the limit.
\end{hint}
\begin{solution}
The $L^\infty$ norm ignores exceptional values on null sets and records the least almost-everywhere uniform bound. As an immediate consequence, the same principle can be inserted into later convergence, approximation, transform, or weak-form arguments whenever its hypotheses are verified.
\end{solution}

\begin{exercise}\label{exr:iii06-04}
State and apply the central principle behind Vector-space structure. Give one immediate consequence in the setting of this chapter.
\end{exercise}
\begin{hint}
On finite-measure spaces, apply Hölder to compare \(L^q\) and \(L^p\) norms when \(q>p\).
\end{hint}
\begin{solution}
Hölder and Minkowski control sums and scalar multiples, making $L^p$ a normed vector space. As an immediate consequence, the same principle can be inserted into later convergence, approximation, transform, or weak-form arguments whenever its hypotheses are verified.
\end{solution}

\begin{exercise}\label{exr:iii06-05}
State and apply the central principle behind Completeness. Give one immediate consequence in the setting of this chapter.
\end{exercise}
\begin{hint}
On infinite-measure spaces, test inclusions with power-law examples instead of assuming the finite-measure relation survives.
\end{hint}
\begin{solution}
Every $L^p$ Cauchy sequence converges in $L^p$ for $1\le p\le\infty$. As an immediate consequence, the same principle can be inserted into later convergence, approximation, transform, or weak-form arguments whenever its hypotheses are verified.
\end{solution}

\begin{exercise}\label{exr:iii06-06}
State and apply the central principle behind Finite-measure inclusions. Give one immediate consequence in the setting of this chapter.
\end{exercise}
\begin{hint}
Use truncation in size and support to approximate an \(L^p\) function by bounded finite-support functions.
\end{hint}
\begin{solution}
On finite-measure spaces, $L^q\subset L^p$ when $q>p$, with an explicit measure factor. As an immediate consequence, the same principle can be inserted into later convergence, approximation, transform, or weak-form arguments whenever its hypotheses are verified.
\end{solution}

\begin{exercise}\label{exr:iii06-07}
State and apply the central principle behind Convergence in measure. Give one immediate consequence in the setting of this chapter.
\end{exercise}
\begin{hint}
Norm convergence does not imply pointwise convergence of the whole sequence; look for an almost-everywhere convergent subsequence if needed.
\end{hint}
\begin{solution}
For finite $p$, $L^p$ convergence implies convergence in measure by Chebyshev's inequality. As an immediate consequence, the same principle can be inserted into later convergence, approximation, transform, or weak-form arguments whenever its hypotheses are verified.
\end{solution}

\begin{exercise}\label{exr:iii06-08}
State and apply the central principle behind Approximation. Give one immediate consequence in the setting of this chapter.
\end{exercise}
\begin{hint}
Check the endpoint \(p=\infty\) separately: the essential supremum replaces an integral formula.
\end{hint}
\begin{solution}
Truncation, finite-measure cutoffs, and simple functions provide dense approximations in $L^p$ for $p<\infty$. As an immediate consequence, the same principle can be inserted into later convergence, approximation, transform, or weak-form arguments whenever its hypotheses are verified.
\end{solution}


% BEGIN VOL03-EXPANSION III06-exercises-01
\section{Graded supplementary exercises}
These exercises extend the preserved foundational set and are graded by mathematical role.

\subsection*{Standard computations}

\begin{exercise}[Power singularity]\label{exr:iii06-expand-01}
For which p is x^(-1/3) in Lp(0,1)?
\end{exercise}
\begin{hint}
Require p/3<1.
\end{hint}
\begin{solution}
Exactly p<3.
\end{solution}

\begin{exercise}[Constant norm]\label{exr:iii06-expand-02}
What is the Lp norm of 1 on a set of measure M?
\end{exercise}
\begin{hint}
Integrate 1.
\end{hint}
\begin{solution}
M^(1/p) for finite p.
\end{solution}

\begin{exercise}[Essential sup]\label{exr:iii06-expand-03}
A function is 7 at one point and 0 elsewhere on [0,1]. What is its Linfinity norm?
\end{exercise}
\begin{hint}
Ignore null exceptions.
\end{hint}
\begin{solution}
It is 0.
\end{solution}

\begin{exercise}[Finite inclusion]\label{exr:iii06-expand-04}
On [0,1], what follows from f in L4?
\end{exercise}
\begin{hint}
Use finite-measure inclusion.
\end{hint}
\begin{solution}
f is in L2 and its L2 norm is at most its L4 norm.
\end{solution}

\begin{exercise}[Chebyshev]\label{exr:iii06-expand-05}
If the Lp norm of f is 2, bound the measure of |f|>t.
\end{exercise}
\begin{hint}
Use Chebyshev.
\end{hint}
\begin{solution}
At most 2^p/t^p.
\end{solution}

\subsection*{Proofs}

\begin{exercise}[Norm zero]\label{exr:iii06-expand-06}
Prove zero Lp norm means zero almost everywhere.
\end{exercise}
\begin{hint}
Apply the zero-integral criterion to |f|^p.
\end{hint}
\begin{solution}
The p-th power is zero almost everywhere.
\end{solution}

\begin{exercise}[Finite inclusion]\label{exr:iii06-expand-07}
Prove Lq is contained in Lp on finite measure when q>p.
\end{exercise}
\begin{hint}
Apply Holder to |f|^p times 1.
\end{hint}
\begin{solution}
The exponent computation gives the stated measure factor.
\end{solution}

\begin{exercise}[Lp to measure]\label{exr:iii06-expand-08}
Prove Lp convergence implies convergence in measure.
\end{exercise}
\begin{hint}
Use Chebyshev on the difference.
\end{hint}
\begin{solution}
Bad-set measure is bounded by norm error to the p divided by epsilon^p.
\end{solution}

\begin{exercise}[Truncation]\label{exr:iii06-expand-09}
Prove bounded truncations converge in Lp.
\end{exercise}
\begin{hint}
Use DCT on the p-th power of the error.
\end{hint}
\begin{solution}
The tail indicator tends to zero under the integrable dominator |f|^p.
\end{solution}

\subsection*{Counterexamples and hypothesis testing}

\begin{exercise}[Infinite measure]\label{exr:iii06-expand-10}
Show finite-measure Lq-to-Lp inclusion can fail on an infinite-measure domain.
\end{exercise}
\begin{hint}
Use power decay on [1,infinity).
\end{hint}
\begin{solution}
Choose an exponent integrable for one power but not the other.
\end{solution}

\begin{exercise}[Pointwise not norm]\label{exr:iii06-expand-11}
Give pointwise convergence without L1 convergence.
\end{exercise}
\begin{hint}
Use moving spikes.
\end{hint}
\begin{solution}
n on (0,1/n) tends to zero almost everywhere but has L1 norm 1.
\end{solution}

\begin{exercise}[Norm not pointwise]\label{exr:iii06-expand-12}
Can Lp convergence fail to imply pointwise convergence of the full sequence?
\end{exercise}
\begin{hint}
Use a typewriter sequence of moving supports.
\end{hint}
\begin{solution}
Yes; norm can tend to zero while points are revisited infinitely often.
\end{solution}

\subsection*{Applications and investigations}

\begin{exercise}[Energy]\label{exr:iii06-expand-13}
Why is L2 natural for signal energy?
\end{exercise}
\begin{hint}
Squared magnitude and orthogonality behave well.
\end{hint}
\begin{solution}
L2 is a Hilbert space and Plancherel preserves its norm.
\end{solution}

\begin{exercise}[Worst case]\label{exr:iii06-expand-14}
Why does Linfinity encode worst-case magnitude only almost everywhere?
\end{exercise}
\begin{hint}
It uses essential supremum.
\end{hint}
\begin{solution}
Null spikes are ignored.
\end{solution}

\subsection*{Challenge problems}

\begin{exercise}[L1-Linfinity interpolation]\label{exr:iii06-expand-15}
If f is in L1 and Linfinity, prove ||f||_p^p <= ||f||_infinity^(p-1)||f||_1.
\end{exercise}
\begin{hint}
Factor |f|^p.
\end{hint}
\begin{solution}
Bound p-1 factors by the essential supremum and integrate.
\end{solution}

\begin{exercise}[Completeness strategy]\label{exr:iii06-expand-16}
Outline the subsequence proof that Lp is complete.
\end{exercise}
\begin{hint}
Choose a subsequence with summable norm increments.
\end{hint}
\begin{solution}
The telescoping increments converge almost everywhere and in norm to a candidate limit.
\end{solution}
% END VOL03-EXPANSION III06-exercises-01
\section*{Chapter summary}
The chapter now supplies a canonical theorem-and-dossier layer for subsequent measure, Fourier, distributional, Sobolev, and PDE arguments.
